% Colonne montante R+10 : pression décroissante, réducteurs aux étages RDC à R+4
\begin{tikzpicture}[font=\scriptsize, >=Stealth, y=0.48cm]
  \draw[blue!70!black, line width=2pt] (0,0) -- (0,10.6);
  \node[draw=vertnuit, fill=vertclair, rounded corners=2pt, align=center] at (0,-1.1)
       {Surpresseur\\ \SI{4,3}{bar}};
  \foreach \k/\p in {0/4{,}3,1/4{,}0,2/3{,}7,3/3{,}4,4/3{,}1,5/2{,}8,6/2{,}5,7/2{,}2,8/1{,}9,9/1{,}6,10/1{,}3} {
    \draw[gray!40] (-2.6,\k) -- (6.2,\k);
    \pgfmathtruncatemacro{\red}{\k<5 ? 1 : 0}
    \ifnum\k=0 \node[anchor=east, gray] at (-1.4,\k+0.35) {RDC};
    \else \node[anchor=east, gray] at (-1.4,\k+0.35) {R+\k}; \fi
    \node[anchor=east] at (-0.15,\k+0.35) {\p~bar};
    \ifnum\red=1
      \draw[blue!70!black, line width=1pt] (0,\k+0.35) -- (1.2,\k+0.35);
      \node[draw=acompleter, fill=acompleter!15, minimum size=3.2mm, inner sep=0pt] at (1.4,\k+0.35) {};
      \draw[blue!70!black, line width=1pt,->] (1.6,\k+0.35) -- (3.0,\k+0.35)
           node[right, black] {3,0 bar (réglé)};
    \else
      \draw[blue!70!black, line width=1pt,->] (0,\k+0.35) -- (3.0,\k+0.35)
           node[right, black] {\p~bar};
    \fi
  }
  \node[draw=acompleter, fill=acompleter!15, minimum size=3.2mm, inner sep=0pt] at (-2.2,-3.2) {};
  \node[anchor=west] at (-2.0,-3.2) {Réducteur de pression + manomètre (étages où $P > 3$~bar)};
\end{tikzpicture}
